In this section, we show bounds on how much the expected discounted utility can be changed by self-modifications of $\epsilon$-optimizers.
\begin{repeatthm}{thm:policy_modification}
Let $\rho$ and $u$ be modification-independent. Consider a self-modifying agent which is $\epsilon$-optimizer for the empty history. Then, for every $t \geq 1$, \begin{gather}
E_{\mae_{<t}} [Q\left(\mae_{<t} \pi_{t}\left(\mae_{<t}\right)\right)] \geq E_{\mae_{<t}}[Q\left(\mae_{<t} \pi_{1}\left(\mae_{<t}\right)\right)] - \min(\frac{\epsilon}{\gamma^{t-1}}, \frac{1}{1-\gamma}) \label{eq:1upper_bound_policy_mod}
\end{gather}
where the expectation is with respect to $\mae_{<t}$ such that the perceptions are distributed according to the belief and the actions are given by $a_{i}=\pi_{i}\left(\mae_{<i}\right)$.

Moreover, for all histories $\mae_{<t}$ given by $a_{i}=\pi_{i}\left(\mae_{<i}\right)$ for which the agent is an $\epsilon$-optimizer it holds that
\begin{gather}
Q\left(\mae_{<t} \pi_{1}\left(\mae_{<t}\right)\right) \geq Q\left(\mae_{<t} \pi_{t}\left(\mae_{<t}\right)\right) - \epsilon \label{eq:1lower_bound_policy_mod}\\
\end{gather}
Equality in Inequality \eqref{eq:1upper_bound_policy_mod} can be achieved up to a factor of at most $\gamma$.
\end{repeatthm}
\begin{proof}[Proof:]
\item
\vspace{-15px}
\paragraph{Inequality~\eqref{eq:1upper_bound_policy_mod}} If the future discounted utility lost, in expectation with respect to $\mae_{<t}$, at time $t$ by $\pi_t$ is $\delta$, then the discounted utility lost at time $1$ is at least $\gamma^{t-1} \delta$. Because the agent is an $\epsilon$-optimizer for the empty history, it holds that  $\gamma^{t-1} \delta \leq \epsilon$. We, therefore, have $\delta \leq \frac{\epsilon}{\gamma^{t-1}}$. To see the other branch of the $\min$, note that no more than the maximum achiavable utility $\frac{1}{1-\gamma}$ can be lost at any time step.

\paragraph{Inequality \eqref{eq:1lower_bound_policy_mod}} follows directly from the definition of an $\epsilon$-optimizer -- the $Q$-value at the histories for which the agent is an $\epsilon$-optimizer has to be at most $\epsilon$ lower than the maximum achiavable one and thus also than $Q\left(\mae_{<t} \pi_{1}\left(\mae_{<t}\right)\right)$. This holds because, thanks to modification-independence, the maximum achievable utility does not depend on the specific policy name and its changes.

\paragraph{Tightness of bound~\eqref{eq:1upper_bound_policy_mod}} can be shown by constructing such a utility function, belief and policy. Let us consider the trivial single-perception deterministic belief, set of two actions $\{0,1\}$, utility function which assigns utility $0$ to histories ending with $0$ and utility $1$ to histories ending with $1$. Now consider the family of policies $\{\pi_i\}_{i=1}^\infty$ such that $\pi_i$ always selects $\pi_{i+1}$ as the next policy and performs action $1$ for $i \leq \frac{\ln(1-\gamma)\epsilon}{\ln \gamma} + 1 \eqdef \mathfrak{b}$
and action $0$ for $i > \mathfrak{b}$. It is easy to verify that the discounted utility lost is equal to at most $\gamma^{\mathfrak{b}-1}\frac{1}{1-\gamma} = \epsilon$, meaning that the policy $\pi_1$ is $\epsilon$-optimizing.

It is also easy to check that the discounted utility lost at time $t$ is
\begin{gather}
\min(\gamma^{\ceil{\mathfrak{b}}-t},1)\frac{1}{1-\gamma} \geq \min(\gamma^{\mathfrak{b}+1-t},1)\frac{1}{1-\gamma} = \gamma^{\mathfrak{b}-1} \frac{1}{1-\gamma} \min(\gamma^{2-t},\gamma^{1-\mathfrak{b}}) \geq \\
\geq \epsilon \gamma \min(\gamma^{1-t},\gamma^{-\mathfrak{b}}) = \epsilon \gamma \min(\gamma^{1-t},\frac{\gamma^{-1}}{(1-\gamma)\epsilon}) = \gamma \min(\frac{\epsilon}{\gamma^{t-1}}, \frac{\gamma^{-1}}{1-\gamma}) \geq \gamma \min(\frac{\epsilon}{\gamma^{t-1}}, \frac{1}{1-\gamma})
\end{gather}
which proves the theorem. Moreover, we see that if $\frac{1}{1-\gamma}$ is the smaller branch of $\min()$ in the bound, the bound is tight.
\end{proof}

From this, we now recover Theorem 16 from \cite{everitt}.

\begin{corollary} \label{cor:recover_everitt}
Let $\rho$ and $u$ be modification-independent. Consider a self-modifying $\epsilon$-optimizer. Then, for every $t \geq 1$, perception sequence $e_{<t}$, and the action sequence $a_{<t}$ given by $a_{i}=\pi_{i}\left(\mae_{<i}\right)$, it then holds that
\begin{gather}
Q\left(\mae_{<t} \pi_{t}\left(\mae_{<t}\right)\right)] = Q\left(\mae_{<t} \pi_{1}\left(\mae_{<t}\right)\right)
\end{gather}
almost surely.
\end{corollary}
\begin{proof}
By \Cref{thm:policy_modification}, it follows that 
\begin{gather}
Q\left(\mae_{<t} \pi_{1}\left(\mae_{<t}\right)\right) \geq Q\left(\mae_{<t} \pi_{t}\left(\mae_{<t}\right)\right)
\end{gather}
and that 
\begin{gather}
E_{\mae_{<t}}[Q\left(\mae_{<t} \pi_{1}\left(\mae_{<t}\right)\right)] \leq E_{\mae_{<t}} [Q\left(\mae_{<t} \pi_{t}\left(\mae_{<t}\right)\right)] 
\end{gather}

Since we have that one inequality holds and the other holds in expectation, it must be the case that the two values are equal almost everywhere.
\end{proof}

Note that while the statement in \cite{everitt} does says that the statement holds surely and not merely almost surely, it is easy to come up with a counterexample to that statement, showing that the fact that we have shown that the statement holds almost surely is a fundamental property of the problem and not only artifact of our analysis.
